\subsection{Sub-figures Example}
\begin{figure}[H]
    \centering
    \begin{subfigure}[t]{0.450\textwidth}
        \includegraphics[width=\textwidth]{Images/Sub-figures_Example/cavity_100_streamlines.png}
        \caption*{$Re = 100$ - Exercise}
    \end{subfigure}
    \begin{subfigure}[t]{0.441\textwidth}
        \includegraphics[width=\textwidth]{Images/Sub-figures_Example/cavity_100_streamlines_ghia.png}
        \caption*{$Re = 100$ \parencite{ghia_1982}}
    \end{subfigure}
    \\  % This makes a new line
    \begin{subfigure}[t]{0.450\textwidth}
        \includegraphics[width=\textwidth]{Images/Sub-figures_Example/cavity_400_streamlines.png}
        \caption*{$Re = 400$ - Exercise}
    \end{subfigure}
    \begin{subfigure}[t]{0.438\textwidth}
        \includegraphics[width=\textwidth]{Images/Sub-figures_Example/cavity_400_streamlines_ghia.png}
        \caption*{$Re = 400$ \parencite{ghia_1982}}
    \end{subfigure}
    \\  % This makes a new line
    \begin{subfigure}[t]{0.450\textwidth}
        \includegraphics[width=\textwidth]{Images/Sub-figures_Example/cavity_1000_streamlines.png}
        \caption*{$Re = 1000$ - Exercise}
    \end{subfigure}
    \begin{subfigure}[t]{0.438\textwidth}
        \includegraphics[width=\textwidth]{Images/Sub-figures_Example/cavity_1000_streamlines_ghia.png}
        \caption*{$Re = 1000$ \parencite{ghia_1982}}
    \end{subfigure}
    \caption[Streamline results]{Streamlines for the problem of a lid-driven cavity.}
    \label{fig:cavity_streamline}
\end{figure}